% ------------------------------------------------------------
% CHAPTER 25
% ------------------------------------------------------------

\chapter{The Geometry of Error: Curved Manifolds and Delta–Sigma Loops}

In the classical theory of delta–sigma modulation, error is treated primarily as a stochastic sequence whose spectral distribution is shaped by the loop filter. The engineering focus lies on the noise transfer function and the power spectral density of the quantization noise. This treatment, while effective, obscures the deeper structure underlying error dynamics. When delta–sigma modulation is examined through the lens of differential geometry, the error becomes not a random variable but a geometric object: a displacement between two points on a curved manifold. The trajectory of the error is governed by curvature, geodesic flows, and the connections defined by the loop filter. This chapter develops a geometric theory of error, demonstrating that delta–sigma stability and resolution can be understood as consequences of geodesic curvature, sectional curvature, and curvature-shaped vector flows.

Let the signal \( \Phi(t) \) be an element of a differentiable manifold \(M\), equipped with a Riemannian metric \(g\). The internal state of the delta–sigma loop, \( x(t) \), is another point on the same manifold. The error is then a tangent vector at \(x\), defined formally by the Riemannian logarithmic map
\[
e(t) = \mathrm{Log}_{x(t)} \Phi(t),
\]
which lies in the tangent space \(T_{x(t)}M\). This definition interprets error not as a difference in Euclidean space but as the initial direction and magnitude of the geodesic that connects the internal state to the true signal. The loop filter acts on this tangent vector, transporting it through the manifold according to the connection compatible with the metric. The dynamics of the error therefore depend on curvature.

To understand these dynamics, we consider the evolution of the error field along a curve \(x(t)\). Differentiating the logarithmic map yields the covariant derivative
\[
\frac{D e}{d t} = \nabla_{\dot{x}} e,
\]
which measures how error changes as the internal state moves through the manifold. In a flat Euclidean space, this derivative reduces to a standard time derivative. On a curved manifold, it includes additional terms involving the Christoffel symbols. The presence of curvature means that even if the signal \(\Phi\) is constant, the error may evolve due to the intrinsic geometry of the manifold. This phenomenon is precisely what occurs in delta–sigma modulators when the signal evolves on a nonlinear manifold or when the internal dynamics impose geometric constraints.

The loop filter defines a vector field \(v\) on the manifold, representing the direction in which the internal state moves to correct error. In geometric terms, the update rule of the delta–sigma loop can be written as
\[
\dot{x}(t) = v(x(t), e(t)),
\]
where the vector field depends on the tangent error. The error dynamics follow from the differential equation
\[
\frac{D e}{d t} = -J_{v}(x(t)) \cdot e(t) - \mathrm{Curv}(x(t), \dot{x}(t), e(t)),
\]
where \(J_{v}\) is the Jacobian of the vector field and \(\mathrm{Curv}\) denotes the curvature term associated with parallel transport. The curvature term involves the Riemann curvature tensor \(R\) and can be written as
\[
\mathrm{Curv}(x(t), \dot{x}(t), e(t)) = R(e(t), \dot{x}(t))\dot{x}(t).
\]
This term quantifies how the geometry of the manifold introduces additional bending of the error trajectory. In a flat space, this term vanishes. In a curved space, it can either amplify error, dampen it, or redirect it into higher-order modes, depending on the sign and magnitude of the sectional curvature.

These geometric insights provide a powerful explanation for phenomena observed in nonlinear delta–sigma modulators. For example, when the error surface is flat, as in linear systems, the error evolves according to a linear differential equation. This yields the familiar high-frequency noise shaping. However, when the error surface is curved, as in modulators with nonlinear quantizers, pole distortion, or state-dependent loop gains, the curvature term can cause error trajectories to exhibit sensitive dependence on initial conditions. This sensitivity is the geometric source of chaos in delta–sigma modulation. Chaos is not merely a nonlinear phenomenon; it is the manifestation of negative curvature on the error manifold.

Conversely, positive curvature can stabilize the loop. In regions of positive sectional curvature, geodesics tend to converge. Error vectors transported along such regions shrink naturally. This explains why certain delta–sigma architectures remain stable even with aggressive loop filters or high gain. Their internal geometry channels the error into regions of positive curvature, producing natural damping. The loop becomes a curvature-regulated engine that maintains coherence by exploiting the geometry of the error surface.

The Riemannian metric also influences quantization behavior. In Euclidean quantization, the quantizer seeks the nearest neighbor among a set of evenly spaced points. In geometric quantization, the quantizer seeks the symbol \(a \in \mathcal{A}\) that minimizes the Riemannian distance
\[
d(x, a) = \sqrt{g_{ij}(x) (x^{i} - a^{i})(x^{j} - a^{j})}.
\]
The Voronoi regions associated with the quantizer become geodesic cells shaped by curvature. As the internal state moves through \(M\), it crosses these curved boundaries, producing discontinuous symbol transitions. These transitions generate entropy injections whose magnitude depends not only on the quantizer spacing but also on the curvature of the manifold. Thus, curvature shapes the statistics of quantization error.

The geometry of error connects naturally to information geometry. In information geometry, the Fisher metric defines the curvature of a statistical manifold. When the delta–sigma loop performs inference, its error evolves on an information manifold whose curvature encodes the predictability of the signal. High curvature corresponds to rapidly changing statistical structure, which makes error correction more difficult. Low curvature corresponds to stable statistical structure, which is easily tracked by the loop. The Fisher metric therefore determines the ease with which a delta–sigma modulator can reconstruct the underlying field.

This geometric interpretation also sheds light on stability. Classical stability criteria analyze the loop gain, poles, and spectral radius. In the geometric picture, stability requires that error trajectories remain within a bounded region of the tangent bundle. This requires that the geodesic curvature of the error flow be bounded. If the sectional curvature is too negative or the Jacobian of the vector field too large, error explodes. If curvature is positive and the Jacobian is moderate, the loop remains stable. Stability becomes a geometric property, not merely an algebraic one.

The energy landscape interpretation of delta–sigma modulation also fits naturally into the geometric setting. The free-energy functional
\[
\mathcal{F}(x) = \frac{1}{2} \| e \|^{2}
\]
defines a potential surface whose gradient flows dictate the motion of the internal state. When the manifold is curved, the gradient is computed using the Riemannian metric:
\[
\nabla \mathcal{F} = g^{ij} \frac{\partial \mathcal{F}}{\partial x^{j}} \frac{\partial}{\partial x^{i}}.
\]
The loop filter therefore implements a Riemannian gradient descent. This descent is shaped by the curvature of the manifold and by the metric itself, which can vary across states. The delta–sigma loop searches for the geodesic that minimizes error, smoothing out high curvature through the integrator dynamics.

In summary, the delta–sigma modulator becomes a dynamical system whose error evolves on a curved manifold governed by Riemannian geometry. Curvature shapes stability, spectral behavior, and quantization statistics. Geodesic deviation governs error propagation. The loop filter defines a connection that transports error along the manifold. The modulator’s behavior is a geometric phenomenon, and the classical linear theory is merely the special case in which the manifold is flat.

The next chapter develops an even deeper connection between error dynamics and lamphrodynamic relaxation, showing that delta–sigma loops implement entropy-minimizing flows that parallel the RSVP theory of field evolution.


